\chapter{其他路线、有效启发与已经看清的障碍}\label{ch:obstacles}
\section{允许复制以后，G条件已经是完整答案}
有一个很强、却必须保留“复制”二字的定理。它最早在三平均研究中整理，证明对任意元数都适用。\src{S28}

\begin{theorem}[复制稳定化]
对非零本原整数零和状态 $X$，以下三件事等价：
\begin{enumerate}
\item $G(X)$ 的所有素因子都整除 $k$；
\item 存在 $R=k^h$，将每个坐标复制 $R$ 份后，能够有限步 $k$ 平均归零；
\item 存在条目属于 $\Z[1/k]$ 的严格正双随机矩阵 $P$，使 $PX=0$。
\end{enumerate}
\end{theorem}
\begin{proof}
先由差分的贝祖等式找整数列 $a$，使 $\mathbf1^Ta=0$、$X^Ta=G$。例如将 $G$ 写成各 $x_i-x_1$ 的整数线性组合即可。置
\[b=Ge_1-x_1a,
\qquad\mathbf1^Tb=G,\qquad X^Tb=0.
\]
取充分大的 $R=k^h$ 使 $G\mid R$，写 $R=nq+r$、$0\le r<n$。因为 $G\mid n$，故 $G\mid r$。令
\[m=q\mathbf1+(r/G)b,\qquad
\ell=R\mathbf1-(n-1)m.
\]
当 $q$ 足够大，两列的所有项都是正整数。它们各自的坐标和为 $R$，且与 $X$ 内积都为零。

取矩阵 $M$ 的前 $n-1$ 行为 $m^T$、最后一行为 $\ell^T$。每行和、每列和都是 $R$，且 $MX=0$。所以 $P=M/R$ 是第三项的矩阵；它的秩还至多为二。

矩阵 $M$ 也直接给复制实现：每个原坐标现有 $R$ 份副本，第 $i$ 组取 $M_{ij}$ 份第 $j$ 种值。列和保证恰好用完副本，行和保证每组有 $R=k^h$ 项，$MX=0$ 保证每组和为零。每组用标准 $k$ 进网络平均，即得第二项。

反过来，若某个 $\ell\nmid k$ 整除 $G$，原输入模 $\ell$ 全同余且非零。复制后的任何平均仍保留这个共同余数；而 $\Z[1/k]$ 上的双随机矩阵也因行和为一保留它。因此第二、三项都不可能。
\end{proof}

它说明 $G$ 已完整刻画了一个放宽后的运输问题。原问题还要求同样的消去发生在原来的有限位置上。正双随机矩阵虽写得很漂亮，却未必能在原位置分解为全等权的 $k$ 平均。

24平均27元是这一差别的压力测试：原位置有已证的 $G=1$ 反例，充分复制后却一定可解。因此不能普遍消去副本，也不能把任意非负运输表直接当作原位置平均词。

\section{二点T变换为何不能直接代替我们的原子}
经典主序理论将许多双随机作用分解为两点混合
\[\begin{pmatrix}\lambda&1-\lambda\\1-\lambda&\lambda\end{pmatrix},
\qquad0\le\lambda\le1.
\]
它允许连续的参数 $\lambda$。我们的操作每次必须把恰好 $k$ 个位置全部替换成同一个均值，既没有任意两点参数，也不能随意保留这些位置上的差异。

因此若要借主序理论，真正需要补的是一个模拟引理：用现有位置和允许的平均精确实现所需混合，并且能连续串联。第\ref{ch:transfer}讲的等大块均值模拟已经给出一种有效情形；一般加权运输仍缺这个分解。

\section{单点交换有能力，也有全长度限制}
块与单点交换能转移有限模见证，但未必能在有理数上提供正向逆。项目对素数平均元数 $p$ 的单大块、至少两个单点系统证明了一个刚性结论。\src{S27}

其差分矩阵为
\[A_i=I-e_ie_i^T-\frac1p\mathbf1e_i^T.
\]
每个 $A_i^{-1}$ 都是整数矩阵，$v_p(\det A_i)=-1$。对长度 $L$ 的非空词 $W$，有
\[p^LW\equiv(-1)^L\mathbf1e_{i_1}^T\pmod p.
\]
这控制了Smith标准形的赋值：在 $m\ge2$ 个差分坐标上，指数为 $(-L,0,\ldots,0)$。一个非零标量矩阵的全部Smith指数应相同，所以这里不可能出现非空标量周期。

特别地，这套单点交换字母表内不存在某个生成元的正向射影逆，否则将它接在该生成元后面就得到标量词。这是所有长度的障碍，不是某个搜索深度没找到。

该结论的原证明针对素数载体。这里引用它来解释一种受限方法为何会卡住，没有把它未经证明地推广到所有合数。一般任务中仍可加入多单点操作、改变核心，或只在有限模层恢复逆；后者正是平方入口所用的方法。

\section{每次平均核很大，也不保证当前参数维数下降}
一次 $k$ 平均消掉选中位置的 $k-1$ 维差分空间。不过若当前状态族已经落在一个线性子空间 $L$ 中，真正降低的维数是
\[\dim(L\cap\ker T),
\]
不是整个 $\dim\ker T=k-1$。核很大，却可能与当前像只交于零。

这解释了为何“每次多并几个单点，总会压到二维”的直觉尚不足以构成证明。必须找到与当前参数族有非零交的实际消去方向，并保证同余合法性。一般移动见证入口只解决了前半段资源准备。

同理，原子差分格 $A_{k-1}$ 的秩为 $k-1$，不意味着证明必须进入 $\mathrm{PGL}_{k-1}$。三种块值加一个零和关系仍只有两个参数，无论 $k$ 多大。高秩算术几何可以提示结构，但参数维数应从当前实际核心计算。\src{S24}

\section{逆向设计零和组：一个可保留的具体族}
另一条路线是先规定第二步出现零和组，再倒推第一次的参数条件。它不替代一般判据，但能给一整族显式短词。\src{S22}

例如奇数 $k\ge5$，在 $n=2k+4$ 个位置取
\[(a^k,b^k,c^2,d^2),
\]
其中
\[
c=\frac{(k^2-2k+8)a-(k^2-6k+6)b}{2},
\]
\[
d=\frac{(-k^2+k-8)a+(k^2-7k+6)b}{2}.
\]
它总和为零。先平均 $(a^{k-4},b^3,c)$，得到 $k$ 份
\[z=\frac{ka+(6-k)b}{2}.
\]
随后平均 $(d,c,z,b^{k-3})$，其和恰为零。所有位置库存都足够；两步后有 $k$ 个零，非零支持至多 $k+4\le2k$，可接既有相反块终端程序，整族至多七步。

这给出了一个二维输入族，不是该维数所有 $G$ 合法输入的判据。它说明“先设计终端，再找能进入它的族”很有用，但还需要一个覆盖全部输入的入口。

同一研究中，两个不交碰撞组的保护与合并也提供了通用组合工具：若一次混合撞上旧重值，另一次可保留关键位置，使两个重值的库存同时增加。非单位差分的中点例外则已由第\ref{ch:tail}讲处理。这些局部工具可复用，不等于区间接续已经完成。

\section{哪些旧工作被后续结果替代了}
旧的36个CRT水平曾逐一编译端点群回路，留下32112条精确回路核验。它们说明特定水平能实现，也保留了短词数据。如今全部 $2k+1$ 端点已有一般证明，不能再把这36个水平以外都列为端点未解；旧证书的具体构造与计算记录仍保留。\src{S23}

另一份早期报告依据一篇clique gossip论文的摘要，把“$k\mid n$”作为固定网络必要条件。六平均九元的实际四步网络说明这句话不能直接用于本模型。项目未取得该论文完整定理及其全部假设，故撤回的是这项本地引用推论，并非声称外部论文错误。\src{S05}

还有两种需要保留的研究边界。第一，某个固定模板的周期失败不等于原问题失败；第二，某个有限生成群不够大也不等于所有合法状态不可达。可以有状态相关的路径，或换一个核心。后续研究应瞄准这些具体接口，不再把历史中某一个充分条件当成原问题的唯一出口。
